% Innioasis Y1 -- manual install via mtkclient (BROM mode).
% The Y1 is a hosted-Linux target: installing Rockbox replaces the original
% firmware on the BOOTIMG and SYSTEM partitions.  There is no Rockbox Utility
% support and nothing is copied onto a mounted drive; both images are written
% with mtkclient while the player is in its boot ROM (BROM) mode.

Installing Rockbox on the \playertype{} replaces the original firmware. It is
written with \fname{mtkclient}, a free flashing tool that talks to the
\daps{} boot ROM (BROM), so no part of the original firmware needs to be
running. The download contains two image files:

\begin{description}
\item[The boot image.] A small image written to the \fname{BOOTIMG} partition.
  It contains the Linux kernel and a minimal start-up program that locates and
  launches Rockbox. It rarely changes and only needs to be written once.
\item[The system image.] An image written to the \fname{ANDROID} partition. It
  is a filesystem holding the main Rockbox binary and the \fname{.rockbox}
  directory (codecs, plugins, fonts and themes).
\end{description}

\warn{Flashing replaces the original firmware entirely; the \playertype{} has
no dual boot. Before you start, make a backup of the stock \fname{boot.img} and
\fname{system.img} so you can return to the original firmware later (see
\reference{sec:uninstall}). Only ever write the \fname{BOOTIMG} and
\fname{ANDROID} partitions -- writing any other partition (in particular the
preloader) can leave the \dap{} unrecoverable without specialist tools.}

\begin{enumerate}
\item Install \fname{mtkclient} on your computer by following the instructions
  at \url{https://github.com/bkerler/mtkclient}. This provides the \fname{mtk}
  command used below.

\item Download your chosen version of Rockbox from the links in the previous
  section and unzip it. You will have two image files, one for the boot
  partition and one for the system partition.

\item Power the \dap{} off completely and disconnect it from the computer.

\item Put the \dap{} into BROM mode and connect it to the computer with a USB
  cable. On the \playertype{} this is done by holding the recessed reset
  button (reachable through the pin-hole) while plugging in the cable.

\item Write both images, substituting the file names of the images you
  downloaded:
\begin{code}
    mtk w bootimg <boot-image-file>
    mtk w android <system-image-file>
\end{code}
\end{enumerate}
